\section{Retrieval}
\label{sec:rag}
\textbf{Pipeline} (\texttt{rag/}). Chunk clauses, embed them with \texttt{bge-small-en-v1.5} (384-d) through the
course embedding endpoint, and index them. A search takes the BM25 top 30 and the dense top 30, fuses them with
reciprocal rank fusion (RRF, $k{=}60$), reranks the pool with a local cross-encoder
(\texttt{Xenova/ms-marco-MiniLM-L-6-v2}, 22M parameters, CPU) and returns the top $k$. The course offers no
reranking endpoint, and using Qwen3-8B as a reranker would cost one extra call per search on the shared
one-request key. A local copy of the same bge weights (fastembed ONNX; cosine 1.00000 against the endpoint) is a
query-time fallback when the endpoint is busy. Every chunk carries a contextual header \texttt{[citation] heading}
and its \texttt{clause\_ids}, which are the basis for citations.

\textbf{Chunking experiments.} \texttt{rag/chunking.py} implements \emph{clause-pack} (whole clauses from one
section packed up to a token budget of 200, 256 or 350, with header) and fixed \emph{token} windows (256 tokens,
32 overlap), each with and without the contextual header. On the corpus this gives 3,881 (clause-pack 200), 3,034
(256), 2,387 (350) and 2,744 (token 256/32) chunks.
\todo[Rajesh Lalam]{chunk-size statistics per strategy (median / p95 / max tokens).}

\textbf{Retrieval recall, measured separately from answers} (\texttt{rag/experiments.py}). We use three metrics.
\emph{sec@k}: a chunk from the gold section is in the top $k$. \emph{cl@k}: the share of gold clauses found.
\emph{fact@k}: the top-$k$ text contains every required fact. fact@k needs no labels and also credits the same rule
found in another source, for example a statute vs.\ the OPM fact sheet. Table~\ref{tab:retr} shows the dev set.

\begin{table}[h]
\centering\small
\caption{Retrieval on the 37 dev-set questions that have gold clauses (6 chunkings $\times$ 4 retrievers; selected
rows).}
\label{tab:retr}
\begin{tabular}{@{}llcccccr@{}}
\toprule
Chunking & Retriever & sec@1 & sec@5 & cl@5 & fact@5 & MRR & ms/q \\
\midrule
clause-pack 200 & hybrid+rerank & 0.189 & \textbf{0.622} & 0.541 & \textbf{0.946} & \textbf{0.405} & 436 \\
clause-pack 256 & hybrid+rerank & 0.216 & 0.595 & 0.541 & \textbf{0.946} & 0.377 & 509 \\
clause-pack 350 & hybrid+rerank & 0.135 & 0.595 & 0.568 & 0.919 & 0.327 & 751 \\
clause-pack 350, no header & hybrid+rerank & 0.135 & 0.486 & 0.486 & 0.892 & 0.289 & 680 \\
token 256/32 & hybrid+rerank & 0.189 & 0.514 & 0.514 & 0.892 & 0.330 & 489 \\
token 256/32 + header & hybrid+rerank & \textbf{0.270} & 0.595 & \textbf{0.595} & 0.919 & 0.404 & 598 \\
\midrule
clause-pack 200 & dense & 0.189 & 0.595 & 0.459 & 0.892 & 0.365 & 50 \\
clause-pack 200 & hybrid (RRF) & 0.162 & 0.541 & 0.446 & 0.838 & 0.316 & 50 \\
clause-pack 200 & BM25 & 0.081 & 0.324 & 0.270 & 0.730 & 0.205 & 5 \\
clause-pack 350, no header & dense & 0.054 & 0.351 & 0.324 & 0.838 & 0.156 & 39 \\
\bottomrule
\end{tabular}
\end{table}

\textbf{Findings.}
(1)~\emph{The reranker is the biggest lever.} Over hybrid RRF it adds 8--16 points of sec@5 on every chunking
(e.g.\ token 256 + header 0.432$\to$0.595), at 0.4--0.75\,s per query on CPU against about 50\,ms without it.
(2)~\emph{Contextual headers help.} Clause-pack 350 goes from 0.486 to 0.595 (rerank) and from 0.351 to 0.514
(dense) with the header, for a few tokens per chunk.
(3)~\emph{Equal-weight RRF is often worse than dense alone.} BM25 alone is weak (sec@5 0.32--0.38), and on 3 of 6
chunkings it drags the fused ranking below dense. Yet hybrid+rerank beats dense: BM25 helps as a \emph{candidate
generator} for the reranker, not as an equal voter.
(4)~\emph{fact@5 is far above cl@5} (0.89--0.95 vs.\ 0.49--0.60). The required fact is usually retrieved, often from
a different source than the labelled clause.
The service runs clause-pack 200 + hybrid + rerank (five chunks, about 1k tokens): it is best by sec@5, fact@5 and
MRR.

\textbf{Recall vs.\ end-to-end.} The eval harness records, per question, whether a gold clause was among the chunks
the agent actually retrieved, next to answer correctness. This separates retrieval misses from generation errors.
\todo{agent retrieval recall vs.\ end-to-end accuracy on dev and held-out sets, with the
explanation of any gap.}
\todo[Rajesh Lalam]{run \texttt{rag/experiments.py} on \texttt{heldout\_v1} (chunking $\times$ retriever grid) and add
it to Table~\ref{tab:retr}.}
